\documentclass[11pt]{article}
\usepackage[T1]{fontenc}
\usepackage{amsmath,amssymb,amsthm}
\usepackage[margin=2.5cm]{geometry}
\usepackage[round,authoryear]{natbib}
\usepackage{booktabs}
\usepackage[hidelinks]{hyperref}
\providecommand{\doi}[1]{doi:\href{https://doi.org/#1}{\nolinkurl{#1}}}

\newtheorem{theorem}{Theorem}
\newtheorem{proposition}[theorem]{Proposition}
\newtheorem{lemma}[theorem]{Lemma}
\newtheorem{corollary}[theorem]{Corollary}
\theoremstyle{remark}
\newtheorem{remark}[theorem]{Remark}

\newcommand{\E}{\mathbb{E}}
\newcommand{\Sig}{S}
\newcommand{\wy}{\mathtt y}
\newcommand{\wt}{\mathtt t}

\title{Truncation error of the signature representation\\ of the Kalman--Bucy filter}
\author{Tamearon Mullings}
\date{}

\begin{document}
\maketitle

\begin{abstract}
For a scalar Ornstein--Uhlenbeck signal observed in continuous time, the Kalman--Bucy filter
mean is a linear functional of the signature of the time-augmented observation path
\citep{cohen2023}. Bounds on the error made by truncating this representation are left open
there. In the constant-gain case we prove two-sided bounds for two choices of coefficients.
With the exact series coefficients the error decays like $(\lambda T)^M/M!$ in the truncation
level $M$. With coefficients chosen optimally at each level it is governed by the best $L^2$
polynomial approximation of the filter kernel, and decays like $(\lambda T/4)^M/M!$ up to a
factor polynomial in $M$. Both results agree with simulations, and a Bene\v{s} example
shows that the picture changes for nonlinear filters.
\end{abstract}

% ======================================================================
\section{Introduction}

The signature of a path, the collection of its iterated integrals, gives a linear basis for
functionals of the path \citep{lyons1998}. \citet{cohen2023} show that the Kalman--Bucy filter
\citep{kalman1961} is a linear function of the initial estimate and of the signature of the
time-augmented observation path, so that linear regression on signature features can recover
the optimal linear filter. In practice the signature is truncated at a finite level $M$, and
\citet[App.~H.3]{cohen2023} note that bounding the resulting approximation error is an open
problem. This note gives explicit, two-sided bounds in the simplest setting: a scalar
Ornstein--Uhlenbeck (OU) signal with the filter started at its steady state.

% ======================================================================
\section{Model and notation}

Fix $T>0$ and let $W,V$ be independent standard Brownian motions. The signal and observation are
\begin{align}
  dX_t &= -aX_t\,dt + \sigma\,dW_t, & X_0 &\sim \mathcal N(0,P_\infty),\label{eq:signal}\\
  dY_t &= hX_t\,dt + \gamma\,dV_t, & Y_0 &= 0,\label{eq:obs}
\end{align}
with $a,\sigma,h,\gamma>0$ and $X_0$ independent of $(W,V)$. Here $P_\infty$ is the positive root of the
algebraic Riccati equation
\begin{equation}\label{eq:are}
  \sigma^2 - 2aP - \tfrac{h^2}{\gamma^2}P^2 = 0,
  \qquad\text{so}\qquad
  P_\infty = \frac{\gamma^2}{h^2}\Big(\!-a+\sqrt{a^2+h^2\sigma^2/\gamma^2}\Big).
\end{equation}
We write
\[
  K=\frac{P_\infty h}{\gamma^2},\qquad \lambda = a+Kh=\sqrt{a^2+h^2\sigma^2/\gamma^2},\qquad
  k(u)=Ke^{-\lambda u},\qquad c_X=\max\Big\{P_\infty,\ \frac{\sigma^2}{2a}\Big\}.
\]
The second expression for $\lambda$ follows from \eqref{eq:are}.
For $f\in L^2[0,T]$ let $\|f\|^2=\int_0^T f(u)^2\,du$, let
$I(f)=\int_0^T f(T-s)\,dY_s$, let $\mathcal P_n$ be the polynomials of degree at most $n$, and let
$E_n(f)=\min_{p\in\mathcal P_n}\|f-p\|$.

\paragraph{Signatures.} The time-augmented path is $\hat Y_t=(t,Y_t)$, with letters $\wt$ and $\wy$.
The signature coordinate of a word $w=w_1\cdots w_n$ on $[0,s]$ is defined recursively by
$\Sig^{\varnothing}_{0,s}=1$ and
$\Sig^{w_1\cdots w_n}_{0,s}=\int_0^s \Sig^{w_1\cdots w_{n-1}}_{0,u}\,d\hat Y^{w_n}_u$
\citep{lyons1998}. Only the words $\wy\,\wt^{\,j}$ (one $\wy$ followed by $j$ letters $\wt$) are needed.
For these, the only integral against $Y$ has integrand $1$, so It\^o and Stratonovich
conventions coincide.

% ======================================================================
\section{The signature representation}

\begin{lemma}\label{lem:words}
For $j\ge0$ and $s\in[0,T]$,
$\displaystyle \Sig^{\wy\wt^{j}}_{0,s}=\int_0^s\frac{(s-u)^j}{j!}\,dY_u$. In particular
$\Sig^{\wy\wt^{j}}_{0,T}=I\big(u^j/j!\big)$.
\end{lemma}

\begin{proof}
For $j=0$ both sides equal $Y_s$. For $j\ge1$ the recursion gives
$\Sig^{\wy\wt^{j}}_{0,s}=\int_0^s \Sig^{\wy\wt^{j-1}}_{0,u}\,du$, a pathwise Riemann integral. Iterating
it with Cauchy's formula for repeated integration gives
$\Sig^{\wy\wt^{j}}_{0,s}=\int_0^s \frac{(s-u)^{j-1}}{(j-1)!}\,Y_u\,du$. Put
$\phi(u)=-(s-u)^j/j!$, so $\phi'(u)=(s-u)^{j-1}/(j-1)!$. Since $\phi$ is deterministic and $C^1$,
integration by parts \citep[Ch.~II]{protter2005} gives
\[
  \int_0^s \phi'(u)Y_u\,du=\phi(s)Y_s-\phi(0)Y_0-\int_0^s\phi(u)\,dY_u=\int_0^s\frac{(s-u)^j}{j!}\,dY_u,
\]
using $\phi(s)=0$ and $Y_0=0$.
\end{proof}

\begin{lemma}[Norm equivalence]\label{lem:norm}
For every $f\in L^2[0,T]$,
\[
  \gamma^2\|f\|^2\;\le\;\E\big[I(f)^2\big]\;\le\;\big(\gamma^2+h^2Tc_X\big)\|f\|^2 .
\]
\end{lemma}

\begin{proof}
Let $g(s)=f(T-s)$, so $\int_0^Tg^2=\|f\|^2$. By \eqref{eq:obs}, $I(f)=hA+\gamma B$ with
$A=\int_0^T g(s)X_s\,ds$ and $B=\int_0^T g(s)\,dV_s$. $A$ is a functional of $(X_0,W)$ and $B$ of $V$,
so they are independent. Since $\E B=0$, $\E[AB]=0$. The It\^o isometry gives $\E B^2=\|f\|^2$,
hence $\E I(f)^2=h^2\E A^2+\gamma^2\|f\|^2$, which proves the lower bound.
From \eqref{eq:signal}, $X_t=e^{-at}X_0+\sigma\int_0^te^{-a(t-r)}dW_r$, so
\[
  \E X_t^2=e^{-2at}P_\infty+\big(1-e^{-2at}\big)\frac{\sigma^2}{2a}\le c_X ,
\]
a convex combination of $P_\infty$ and $\sigma^2/2a$. Minkowski's integral inequality and
Cauchy--Schwarz then give
$(\E A^2)^{1/2}\le\int_0^T|g(s)|(\E X_s^2)^{1/2}ds\le\sqrt{c_X}\int_0^T|g|\le\sqrt{c_XT}\,\|f\|$,
which proves the upper bound.
\end{proof}

\begin{proposition}[\citealp{cohen2023}, Thm.~1, constant-gain case]\label{prop:rep}
$\displaystyle m_T:=\E\big[X_T\mid\mathcal F^Y_T\big]=I(k)=\sum_{j\ge0}K(-\lambda)^j\,\Sig^{\wy\wt^{j}}_{0,T}$,
with the series converging in $L^2(\mathbb P)$.
\end{proposition}

\begin{proof}
The Kalman--Bucy filter \citep{kalman1961,bain2009} for \eqref{eq:signal}--\eqref{eq:obs} is
$dm_t=-am_t\,dt+P_th\gamma^{-2}(dY_t-hm_t\,dt)$, with $\dot P_t=\sigma^2-2aP_t-h^2\gamma^{-2}P_t^2$,
$m_0=0$ and $P_0=P_\infty$. Since $P_\infty$ is a fixed point of the Riccati equation,
$P_t\equiv P_\infty$ and $dm_t=-\lambda m_t\,dt+K\,dY_t$. Applying It\^o's product rule to
$e^{\lambda t}m_t$ gives $m_T=\int_0^TKe^{-\lambda(T-s)}dY_s=I(k)$.
Let $k_M(u)=K\sum_{j<M}(-\lambda u)^j/j!$. By Lemma~\ref{lem:words} and linearity,
$I(k_M)=\sum_{j<M}K(-\lambda)^j\Sig^{\wy\wt^{j}}_{0,T}$. Taylor's theorem gives
$\sup_{[0,T]}|k-k_M|\le K(\lambda T)^M/M!\to0$, so $\|k-k_M\|\to0$, and by
Lemma~\ref{lem:norm} $\E[(m_T-I(k_M))^2]\to0$.
\end{proof}

% ======================================================================
\section{Truncation with the exact coefficients}

Let $m^{(M)}_T=\sum_{j<M}K(-\lambda)^j\Sig^{\wy\wt^{j}}_{0,T}=I(k_M)$, the representation truncated at
signature level $M$, and $r_M=k-k_M$.

\begin{theorem}\label{thm:taylor}
For every $M\ge1$,
\[
  \E\big[(m_T-m^{(M)}_T)^2\big]\le\big(\gamma^2+h^2Tc_X\big)\,
  \frac{K^2\lambda^{2M}T^{2M+1}}{(2M+1)(M!)^2}.
\]
If moreover $M+1>\lambda T$, then
$\E[(m_T-m^{(M)}_T)^2]\ge\gamma^2\big(1-\tfrac{\lambda T}{M+1}\big)^2\frac{K^2\lambda^{2M}T^{2M+1}}{(2M+1)(M!)^2}$.
\end{theorem}

\begin{proof}
For $x\ge0$, the Lagrange form of the remainder gives
$\big|e^{-x}-\sum_{j<M}(-x)^j/j!\big|=e^{-\xi}x^M/M!\le x^M/M!$ for some $\xi\in(0,x)$. Hence
$|r_M(u)|\le K(\lambda u)^M/M!$, and $\|r_M\|^2\le K^2\lambda^{2M}T^{2M+1}/\big((2M+1)(M!)^2\big)$.
The upper bound follows from Lemma~\ref{lem:norm}, since $m_T-m_T^{(M)}=I(r_M)$.
For the lower bound, write $r_M(u)/K=\sum_{j\ge M}(-1)^j\tau_j$ with $\tau_j=x^j/j!$ and $x=\lambda u$.
If $x<M+1$, then $\tau_{j+1}/\tau_j=x/(j+1)<1$ for all $j\ge M$. The series is then alternating with
decreasing terms, so $|r_M(u)|/K\ge\tau_M-\tau_{M+1}=\frac{x^M}{M!}\big(1-\frac{x}{M+1}\big)$.
For $u\le T$ the last factor is at least $1-\lambda T/(M+1)>0$. Squaring, integrating and applying
the lower half of Lemma~\ref{lem:norm} gives the claim.
\end{proof}

The bounds match up to constants, so the exact-coefficient error decays like $(\lambda T)^M/M!$.
Their ratio between consecutive levels is $\rho^{\mathrm{T}}_M=\big(\lambda T/M\big)^2\,\frac{2M-1}{2M+1}$.

% ======================================================================
\section{Truncation with optimal coefficients}

Let $\mathcal S_M=\operatorname{span}\{\Sig^{\wy\wt^{j}}_{0,T}:j<M\}$ and let $m^{(M,*)}_T$ be the
$L^2(\mathbb P)$-orthogonal projection of $m_T$ onto $\mathcal S_M$. This is the best mean-square
predictor built from these signature terms, and it is the quantity that least-squares
regression on them estimates.

\begin{theorem}\label{thm:opt}
For every $M\ge1$,
\[
  \gamma^2E_{M-1}(k)^2\;\le\;\E\big[(m_T-m^{(M,*)}_T)^2\big]\;\le\;\big(\gamma^2+h^2Tc_X\big)E_{M-1}(k)^2 .
\]
\end{theorem}

\begin{proof}
By Lemma~\ref{lem:words}, $\mathcal S_M=\{I(p):p\in\mathcal P_{M-1}\}$, since the $u^j/j!$ with
$j<M$ span $\mathcal P_{M-1}$. As $I$ is linear,
$\E[(m_T-m_T^{(M,*)})^2]=\min_{p\in\mathcal P_{M-1}}\E[I(k-p)^2]$. Choosing $p$ as the $L^2[0,T]$-best
approximation of $k$ and applying the upper half of Lemma~\ref{lem:norm} gives the upper bound.
For any $p$, the lower half gives $\E I(k-p)^2\ge\gamma^2\|k-p\|^2\ge\gamma^2E_{M-1}(k)^2$.
\end{proof}

It remains to evaluate $E_n(k)$. Let $i_\ell$ denote the modified spherical Bessel function of
the first kind, $i_\ell=\mathsf i^{(1)}_\ell$ in the notation of \citet[\S10.47]{dlmf}, and let
$c=\lambda T/2$.

\begin{lemma}[Best $L^2$ approximation of the kernel]\label{lem:legendre}
For $n\ge0$,
\[
  E_n(k)^2=K^2e^{-2c}\,\frac T2\sum_{\ell>n}2(2\ell+1)\,i_\ell(c)^2 .
\]
Moreover, with $\theta_\ell=2(2\ell+1)\,c^{2\ell}/\big((2\ell+1)!!\big)^2$ and
$q_n=c^2/\big((2n+3)(2n+5)\big)$, if $q_n<1$ then
\[
  \theta_{n+1}\;\le\;\frac{2}{T}\,K^{-2}e^{2c}E_n(k)^2\;\le\;\theta_{n+1}\,\frac{e^{c^2/(2n+5)}}{1-q_n}.
\]
\end{lemma}

\begin{proof}
The substitution $u=T(1-x)/2$ maps $x\in[-1,1]$ onto $u\in[0,T]$, maps $\mathcal P_n$ onto itself,
and gives $k(u)=Ke^{-c}e^{cx}$ and $\|f\|^2=\frac T2\int_{-1}^1 f(u(x))^2dx$. So
$E_n(k)=Ke^{-c}\sqrt{T/2}\,\mathcal E_n$, where $\mathcal E_n$ is the best $L^2[-1,1]$ error of
$e^{cx}$ by $\mathcal P_n$.

Rayleigh's expansion \citep[eq.~10.60.7]{dlmf} with $z=-\mathrm ic$, together with
$\mathsf i^{(1)}_\ell(c)=\mathrm i^{-\ell}\mathsf j_\ell(\mathrm ic)$ \citep[eq.~10.47.12]{dlmf}
and the parity of $\mathsf j_\ell$, gives the Legendre expansion
$e^{cx}=\sum_{\ell\ge0}(2\ell+1)i_\ell(c)P_\ell(x)$. The Legendre polynomials are orthogonal with
$\int_{-1}^1P_\ell^2=2/(2\ell+1)$, so the best approximation in $\mathcal P_n$ is the truncated series
\citep[Ch.~17]{trefethen2019}, and $\mathcal E_n^2=\sum_{\ell>n}2(2\ell+1)i_\ell(c)^2$.

For the bounds, $i_\ell(c)=c^\ell\sum_{r\ge0}\frac{(c^2/2)^r}{r!\,(2\ell+2r+1)!!}$
\citep[eq.~10.53.3]{dlmf}. All terms are positive, and $(2\ell+2r+1)!!\ge(2\ell+1)!!\,(2\ell+3)^r$, so
\[
  \frac{c^\ell}{(2\ell+1)!!}\le i_\ell(c)\le\frac{c^\ell}{(2\ell+1)!!}\,e^{c^2/(2(2\ell+3))}.
\]
The lower bound follows by keeping only $\ell=n+1$. For the upper bound, $\theta_{\ell+1}/\theta_\ell=c^2/((2\ell+1)(2\ell+3))\le q_n$
for $\ell\ge n+1$, so $\sum_{\ell>n}\theta_\ell\le\theta_{n+1}/(1-q_n)$. Also
$e^{c^2/(2\ell+3)}\le e^{c^2/(2n+5)}$ for $\ell\ge n+1$.
\end{proof}

\begin{corollary}\label{cor:rate}
With $n=M-1$, Theorem~\ref{thm:opt} and Lemma~\ref{lem:legendre} give
\[
  \gamma^2K^2e^{-\lambda T}\tfrac T2\,\theta_M\;\le\;\E\big[(m_T-m^{(M,*)}_T)^2\big]\;\le\;
  \big(\gamma^2+h^2Tc_X\big)K^2e^{-\lambda T}\tfrac T2\,\theta_M\,\frac{e^{c^2/(2M+3)}}{1-q_{M-1}} .
\]
Since $(2M+1)!!=(2M+1)!/(2^MM!)$ and $(2M)!\sim4^M(M!)^2/\sqrt{\pi M}$, the root-mean-square
error is of order $(\lambda T/4)^M/M!$ up to a factor polynomial in $M$. Compared with
Theorem~\ref{thm:taylor} this is smaller by a factor of order $4^{-M}$. Consecutive levels shrink
the mean-square error by $\rho^{*}_M=\theta_M/\theta_{M-1}=c^2/\big((2M-1)(2M+1)\big)\approx(\lambda T/4M)^2$.
\end{corollary}

% ======================================================================
\section{Numerical verification}\label{sec:numerics}

\paragraph{Setup.}
\begin{itemize}
  \item Parameters: $a=\sigma=h=1$, $\gamma=0.5$, $T=1$. Hence $P_\infty=0.3090$, $K=1.2361$,
        $\lambda=\sqrt5=2.2361$, $c_X=0.5$, and $\gamma^2+h^2Tc_X=3\gamma^2$.
  \item Simulation: $6000$ paths of \eqref{eq:signal}--\eqref{eq:obs} on an Euler grid with
        $dt=10^{-3}$; $4000$ for fitting, $2000$ for testing.
  \item Signatures: computed with \texttt{esig} on the piecewise-linear interpolation of $Y$, and
        cross-checked against \texttt{pysiglib} (largest difference $2\times10^{-16}$). The target
        is $I(k)$ evaluated on that interpolated path, which the signature series represents exactly.
  \item Checks: the target agrees with a discrete-time Kalman filter to root-mean-square error
        $1.7\times10^{-4}$.
\end{itemize}
Errors and bounds are reported as mean-square error divided by $\operatorname{Var}(m_T)=0.1651$ (test set).
Code and raw output are in the repository (\texttt{d2\_kb\_signature.py},
\texttt{results/d2/summary.json}).

\begin{table}[h]
\centering
\caption{Exact coefficients (Theorem~\ref{thm:taylor}). The lower bound needs $M+1>\lambda T$.}
\label{tab:taylor}
\begin{tabular}{c c c c c c}
\toprule
$M$ & measured & lower bound & upper bound & measured ratio & $\rho^{\mathrm T}_M$ \\
\midrule
1 & $1.87$                & --                    & $1.16\times10^{1}$    & --     & --     \\
2 & $1.65$                & $1.88\times10^{-1}$   & $8.67$                & 0.886  & 0.750  \\
3 & $7.29\times10^{-1}$   & $2.23\times10^{-1}$   & $3.44$                & 0.441  & 0.397  \\
4 & $1.90\times10^{-1}$   & $8.52\times10^{-2}$   & $8.37\times10^{-1}$   & 0.261  & 0.243  \\
5 & $3.26\times10^{-2}$   & $1.80\times10^{-2}$   & $1.37\times10^{-1}$   & 0.172  & 0.164  \\
6 & $3.96\times10^{-3}$   & $2.48\times10^{-3}$   & $1.61\times10^{-2}$   & 0.122  & 0.118  \\
7 & $3.59\times10^{-4}$   & $2.46\times10^{-4}$   & $1.42\times10^{-3}$   & 0.091  & 0.088  \\
8 & $2.52\times10^{-5}$   & $1.85\times10^{-5}$   & $9.81\times10^{-5}$   & 0.070  & 0.069  \\
\bottomrule
\end{tabular}
\end{table}

\begin{table}[h]
\centering
\caption{Optimal coefficients (Theorem~\ref{thm:opt}, with $E_{M-1}(k)$ from Lemma~\ref{lem:legendre}).}
\label{tab:opt}
\begin{tabular}{c c c c c c c}
\toprule
$M$ & measured & lower bound & upper bound & measured ratio & $\frac{E_{M-1}^2}{E_{M-2}^2}$ & $\rho^{*}_M$ \\
\midrule
1 & $1.80\times10^{-1}$  & $1.42\times10^{-1}$  & $4.27\times10^{-1}$  & --     & --     & --     \\
2 & $1.17\times10^{-2}$  & $1.06\times10^{-2}$  & $3.18\times10^{-2}$  & 0.0645 & 0.0745 & 0.0833 \\
3 & $3.72\times10^{-4}$  & $3.59\times10^{-4}$  & $1.08\times10^{-3}$  & 0.0320 & 0.0338 & 0.0357 \\
4 & $7.02\times10^{-6}$  & $6.90\times10^{-6}$  & $2.07\times10^{-5}$  & 0.0189 & 0.0192 & 0.0198 \\
5 & $8.38\times10^{-8}$  & $8.53\times10^{-8}$  & $2.56\times10^{-7}$  & 0.0119 & 0.0124 & 0.0126 \\
6 & $7.31\times10^{-10}$ & $7.34\times10^{-10}$ & $2.20\times10^{-9}$  & 0.0087 & 0.0086 & 0.0087 \\
7 & $4.68\times10^{-12}$ & $4.65\times10^{-12}$ & $1.40\times10^{-11}$ & 0.0064 & 0.0063 & 0.0064 \\
8 & $2.10\times10^{-14}$ & $2.26\times10^{-14}$ & $6.79\times10^{-14}$ & 0.0045 & 0.0049 & 0.0049 \\
\bottomrule
\end{tabular}
\end{table}

\paragraph{Exact coefficients (Table~\ref{tab:taylor}).} Every measured error lies inside the bounds
of Theorem~\ref{thm:taylor}. The level-to-level ratios approach $\rho^{\mathrm T}_M$ (0.070 against
0.069 at $M=8$).

\paragraph{Optimal coefficients (Table~\ref{tab:opt}).} Here the coefficients are fitted by least
squares on the $4000$ training paths. The measured errors sit almost exactly on the lower bound, so
the observation-noise term dominates the constant in Lemma~\ref{lem:norm}. Two entries ($M=5,8$) fall
$1.7\%$ and $7\%$ below the lower bound. The bound concerns the continuous-time population error,
while the table reports a $2000$-path test-set estimate on a time grid; at $M=8$ the relative error
is also near $10^{-14}$, where floating-point effects in the least-squares fit matter. The ratios
match $E_{M-1}^2/E_{M-2}^2$ and approach $\rho^*_M\approx(\lambda T/4M)^2$. At $M=8$ the error is
about $10^{9}$ times smaller than with the exact coefficients.

% ======================================================================
\section{Outlook: a nonlinear filter}\label{sec:outlook}

For the Bene\v{s} model $dX_t=\tanh(X_t)\,dt+dW_t$, $dY_t=X_t\,dt+\gamma\,dV_t$, $X_0=0$
\citep{benes1981}, the filter is finite-dimensional. Since $\tanh'+\tanh^2=1$, Girsanov's theorem
gives the posterior as $\cosh(x)$ times the Gaussian Kalman--Bucy posterior of a Brownian signal.
Its mean is therefore $m_T+P_T\tanh(m_T)$, where $m_T=\int_0^T\sinh(hs/\gamma)\,dY_s/(\gamma\cosh(hT/\gamma))$
and $P_T=(\gamma/h)\tanh(hT/\gamma)$. Again $m_T$ is a single-$\wy$ signature functional.

In simulations with $\gamma=1.5$ and $T=2$ ($6000$ paths; the closed form agrees with a
$20{,}000$-particle filter to root-mean-square error $0.014$ against a target standard deviation of
$1.91$):
\begin{itemize}
  \item Least squares on the words $\wy\wt^{j}$ stops at $1.25\times10^{-2}$ (relative mean-square
        error) from $M=3$ onwards. This is exactly the share of the target's variance that is
        nonlinear in $m_T$.
  \item Ridge regression on the full truncated signature goes below this floor:
        $6.6\times10^{-3}$, $2.6\times10^{-3}$, $1.7\times10^{-3}$ and $7.8\times10^{-4}$ for
        $M=3,\dots,6$. That is roughly a factor of 2 per level. For $M\ge7$ the error is limited by
        the training-set size.
\end{itemize}
Bounding this slower, geometric rate requires words with several $\wy$, where the It\^o and
Stratonovich lifts differ, and polynomial approximation of $\tanh$ composed with a Gaussian
functional. This is left for future work.

\section*{Acknowledgement}
The numerical experiments and the typesetting of this note were carried out with the assistance
of an AI coding assistant (Claude, Anthropic). The derivations were worked through and checked
by the author.

% ======================================================================
\begin{thebibliography}{9}
\bibitem[Bain and Crisan(2009)]{bain2009}
Bain, A. and Crisan, D. (2009) \emph{Fundamentals of Stochastic Filtering}. New York: Springer.
\doi{10.1007/978-0-387-76896-0}.

\bibitem[Bene\v{s}(1981)]{benes1981}
Bene\v{s}, V.E. (1981) `Exact finite-dimensional filters for certain diffusions with nonlinear
drift', \emph{Stochastics}, 5(1--2), pp.~65--92. \doi{10.1080/17442508108833174}.

\bibitem[Cohen et~al.(2023)]{cohen2023}
Cohen, S.N., Mantoan, G., Nesheim, L., de Paula, \'A., Turrell, A. and Yang, L. (2023)
`Nowcasting using regression on signatures', \emph{arXiv preprint} arXiv:2305.10256 [v3].
Available at: \url{https://arxiv.org/abs/2305.10256} (Accessed: 18 October 2026).

\bibitem[Kalman and Bucy(1961)]{kalman1961}
Kalman, R.E. and Bucy, R.S. (1961) `New results in linear filtering and prediction theory',
\emph{Journal of Basic Engineering}, 83(1), pp.~95--108. \doi{10.1115/1.3658902}.

\bibitem[Lyons(1998)]{lyons1998}
Lyons, T.J. (1998) `Differential equations driven by rough signals', \emph{Revista Matem\'atica
Iberoamericana}, 14(2), pp.~215--310. \doi{10.4171/RMI/240}.

\bibitem[Olver et~al.(2026)]{dlmf}
Olver, F.W.J., Olde Daalhuis, A.B., Lozier, D.W., Schneider, B.I., Boisvert, R.F., Clark, C.W.,
Miller, B.R., Saunders, B.V., Cohl, H.S. and McClain, M.A. (eds.) (2026) \emph{NIST Digital Library
of Mathematical Functions}. Release 1.2.8 of 2026-09-15. Available at: \url{https://dlmf.nist.gov/}
(Accessed: 18 October 2026).

\bibitem[Protter(2005)]{protter2005}
Protter, P.E. (2005) \emph{Stochastic Integration and Differential Equations}. 2nd edn. Berlin:
Springer. \doi{10.1007/978-3-662-10061-5}.

\bibitem[Trefethen(2019)]{trefethen2019}
Trefethen, L.N. (2019) \emph{Approximation Theory and Approximation Practice}. Extended edn.
Philadelphia: SIAM. \doi{10.1137/1.9781611975949}.
\end{thebibliography}

\end{document}
